\documentclass[12pt,a4paper]{article}

\usepackage[utf8]{inputenc}
\usepackage[T1]{fontenc}
\usepackage[spanish]{babel}
\usepackage{amsmath,amssymb}
\usepackage{geometry}
\usepackage{booktabs}
\usepackage{array}
\usepackage{parskip}
\usepackage{fancyhdr}
\usepackage{listings}
\usepackage{xcolor}

\geometry{margin=2.5cm, headheight=15pt}

\definecolor{codebg}{RGB}{245,245,245}
\definecolor{keywordcolor}{RGB}{0,100,180}
\definecolor{commentcolor}{RGB}{80,110,80}
\lstset{
  language=Python,
  backgroundcolor=\color{codebg},
  frame=single, rulecolor=\color{gray!40},
  basicstyle=\ttfamily\small,
  keywordstyle=\color{keywordcolor}\bfseries,
  commentstyle=\color{commentcolor}\itshape,
  numbers=left, numberstyle=\tiny\color{gray}, numbersep=8pt,
  breaklines=true, showstringspaces=false, tabsize=4,
}

\pagestyle{fancy}
\fancyhf{}
\rhead{Modelado y Simulaci\'on --- 22/04/2026}
\lhead{Punto 3 --- Docente: Omar C\'aceres}
\cfoot{\thepage}

\title{\textbf{Examen --- Modelado y Simulaci\'on}\\[0.3em]
       {\large Punto 3: Integraci\'on Num\'erica --- Trapecio y Simpson 1/3}}
\author{Juan Bautista Espino}
\date{22 de abril de 2026}

\begin{document}
\maketitle\thispagestyle{fancy}

% ══════════════════════════════════════════════════════════════════════════════
\section*{Enunciado}
% ══════════════════════════════════════════════════════════════════════════════

Se desea aproximar la integral:
\[
  I = \int_0^1 \frac{1}{\sqrt{2}}\, x^{2x}\, dx
\]

\begin{enumerate}
  \item[(a)] Aplicar la regla del \textbf{Trapecio} con $n=4$ y calcular el error de
             truncamiento en $\xi = 0.5$.
  \item[(b)] Aplicar la regla de \textbf{Simpson 1/3} con $n=4$ y calcular el error de
             truncamiento en $\xi = 0.5$.
\end{enumerate}

% ══════════════════════════════════════════════════════════════════════════════
\section*{Configuraci\'on com\'un ($n=4$)}
% ══════════════════════════════════════════════════════════════════════════════

\[
  a = 0, \quad b = 1, \quad n = 4, \quad
  h = \frac{b-a}{n} = \frac{1}{4} = 0.25
\]

La funci\'on a integrar es:
\[
  f(x) = \frac{1}{\sqrt{2}}\, x^{2x} = \frac{1}{\sqrt{2}}\, e^{2x\ln x},
  \qquad f(0) \coloneqq \lim_{x\to 0^+} \frac{x^{2x}}{\sqrt{2}} = \frac{1}{\sqrt{2}}
\]

El l\'imite $\lim_{x\to 0^+} x^{2x} = \lim_{x\to 0^+} e^{2x\ln x} = e^0 = 1$
se aplica por continuidad.

\subsection*{Tabla de nodos y valores}

\begin{center}
\renewcommand{\arraystretch}{1.35}
\begin{tabular}{ccc}
\toprule
$i$ & $x_i$ & $f(x_i) = \dfrac{x_i^{2x_i}}{\sqrt{2}}$ \\[4pt]
\midrule
0 & 0.00 & $\dfrac{1}{\sqrt{2}} \approx 0.707\,107$ \\[6pt]
1 & 0.25 & $\dfrac{(0.25)^{0.5}}{\sqrt{2}} = \dfrac{0.5}{\sqrt{2}} \approx 0.353\,553$ \\[6pt]
2 & 0.50 & $\dfrac{(0.5)^{1}}{\sqrt{2}} = \dfrac{0.5}{\sqrt{2}} \approx 0.353\,553$ \\[6pt]
3 & 0.75 & $\dfrac{(0.75)^{1.5}}{\sqrt{2}} \approx \dfrac{0.649\,519}{\sqrt{2}} \approx 0.459\,279$ \\[6pt]
4 & 1.00 & $\dfrac{1^2}{\sqrt{2}} = \dfrac{1}{\sqrt{2}} \approx 0.707\,107$ \\[4pt]
\bottomrule
\end{tabular}
\end{center}

\textit{Nota:} $(0.25)^{2\times0.25} = (0.25)^{0.5} = 0.5$ y
$(0.5)^{2\times0.5} = (0.5)^1 = 0.5$, de ah\'i que $f(x_1)=f(x_2)$.

% ══════════════════════════════════════════════════════════════════════════════
\section*{Parte (a) --- Regla del Trapecio ($n=4$)}
% ══════════════════════════════════════════════════════════════════════════════

\subsection*{F\'ormula}

\[
  I_t = \frac{h}{2}\Bigl[f(x_0) + 2\sum_{i=1}^{n-1}f(x_i) + f(x_n)\Bigr]
      = \frac{h}{2}\bigl[f_0 + 2f_1 + 2f_2 + 2f_3 + f_4\bigr]
\]

\subsection*{C\'alculo detallado}

\begin{align*}
  I_t &= \frac{0.25}{2}\bigl[0.707\,107
         + 2(0.353\,553)
         + 2(0.353\,553)
         + 2(0.459\,279)
         + 0.707\,107\bigr] \\[4pt]
      &= 0.125 \times \bigl[0.707\,107
         + 0.707\,107
         + 0.707\,107
         + 0.918\,558
         + 0.707\,107\bigr] \\[4pt]
      &= 0.125 \times 3.746\,986
\end{align*}

\[
  \boxed{I_t \approx 0.468\,373}
\]

\subsection*{Error de truncamiento}

La f\'ormula del error para el Trapecio compuesto es:
\[
  E_t = -\frac{(b-a)\,h^2}{12}\,f''(\xi), \qquad \xi \in (0,1)
\]

\subsubsection*{Derivaci\'on de $f''(x)$}

Escribimos $f(x) = \dfrac{1}{\sqrt{2}}\,e^{u(x)}$ con $u(x) = 2x\ln x$.

\paragraph{Primera derivada} (regla de la cadena):
\[
  u'(x) = 2\ln x + 2x\cdot\frac{1}{x} = 2\ln x + 2 = 2(\ln x + 1)
\]
\[
  f'(x) = f(x)\cdot u'(x) = f(x)\cdot 2(\ln x + 1)
\]

\paragraph{Segunda derivada} (producto $f'\cdot u' + f\cdot u''$):
\[
  u''(x) = \frac{d}{dx}\bigl[2\ln x + 2\bigr] = \frac{2}{x}
\]
\begin{align*}
  f''(x) &= f'(x)\cdot 2(\ln x+1) + f(x)\cdot u''(x) \\
          &= f(x)\cdot\bigl[2(\ln x+1)\bigr]^2 + f(x)\cdot\frac{2}{x} \\[4pt]
          &= f(x)\cdot\left[4(\ln x+1)^2 + \frac{2}{x}\right]
\end{align*}

\subsubsection*{Evaluaci\'on en $\xi = 0.5$}

\begin{align*}
  \ln(0.5) &= \ln\!\tfrac{1}{2} = -\ln 2 \approx -0.693\,147 \\
  \ln(0.5)+1 &= 1 - 0.693\,147 = 0.306\,853 \\
  f(0.5) &= \frac{(0.5)^1}{\sqrt{2}} = \frac{0.5}{\sqrt{2}} \approx 0.353\,553
\end{align*}

\begin{align*}
  f''(0.5)
  &= 0.353\,553 \times \Bigl[4\,(0.306\,853)^2 + \frac{2}{0.5}\Bigr] \\
  &= 0.353\,553 \times \bigl[\,4 \times 0.094\,159 + 4\,\bigr] \\
  &= 0.353\,553 \times \bigl[\,0.376\,637 + 4\,\bigr] \\
  &= 0.353\,553 \times 4.376\,637 \\
  &\approx 1.547\,374
\end{align*}

\subsubsection*{C\'alculo final}

\[
  E_t = -\frac{(1)(0.25)^2}{12}\times 1.547\,374
      = -\frac{0.062\,5}{12}\times 1.547\,374
      = -0.005\,208\times 1.547\,374
\]

\[
  \boxed{E_t \approx -0.008\,059}
\]


% ══════════════════════════════════════════════════════════════════════════════
\section*{Parte (b) --- Regla de Simpson 1/3 ($n=4$)}
% ══════════════════════════════════════════════════════════════════════════════

\subsection*{F\'ormula}

Para $n$ par (aqu\'i $n=4$):
\[
  I_s = \frac{h}{3}\Bigl[f_0 + 4f_1 + 2f_2 + 4f_3 + f_4\Bigr]
\]

\subsection*{C\'alculo detallado}

\begin{align*}
  I_s &= \frac{0.25}{3}\bigl[
         \underbrace{0.707\,107}_{f_0}
         + 4\underbrace{(0.353\,553)}_{f_1}
         + 2\underbrace{(0.353\,553)}_{f_2}
         + 4\underbrace{(0.459\,279)}_{f_3}
         + \underbrace{0.707\,107}_{f_4}\bigr] \\[6pt]
      &= \frac{0.25}{3}\bigl[
         0.707\,107
         + 1.414\,214
         + 0.707\,107
         + 1.837\,117
         + 0.707\,107\bigr] \\[4pt]
      &= 0.083\,333 \times 5.372\,651
\end{align*}

\[
  \boxed{I_s \approx 0.447\,721}
\]

\subsection*{Error de truncamiento}

La f\'ormula del error para Simpson 1/3 compuesto es:
\[
  E_s = -\frac{(b-a)\,h^4}{180}\,f^{(4)}(\xi), \qquad \xi \in (0,1)
\]

\subsubsection*{Derivaci\'on sistem\'atica hasta $f^{(4)}(x)$}

Definimos $u(x)=2x\ln x$, de modo que $f(x)=\tfrac{1}{\sqrt{2}}\,e^{u}$.
Las derivadas de $u$ son:

\[
  u'  = 2(\ln x+1), \quad
  u'' = \frac{2}{x}, \quad
  u'''= -\frac{2}{x^2}, \quad
  u^{(4)} = \frac{4}{x^3}
\]

Aplicando la f\'ormula de Fa\`a di Bruno (o diferenciando sucesivamente
$f^{(k)} = \tfrac{d}{dx}f^{(k-1)}$) se obtiene:

\begin{align*}
  f'(x)   &= f\,u' \\
  f''(x)  &= f\!\left[(u')^2 + u''\right]
           = f\!\left[4(\ln x+1)^2+\tfrac{2}{x}\right] \\
  f'''(x) &= f\!\left[(u')^3 + 3u'u'' + u'''\right]
           = f\!\left[8(\ln x+1)^3+\tfrac{12(\ln x+1)}{x}-\tfrac{2}{x^2}\right] \\
  f^{(4)}(x) &= f\!\left[(u')^4+6(u')^2 u''+3(u'')^2+4u'u'''+u^{(4)}\right]
\end{align*}

Desarrollando el \'ultimo t\'ermino:
\begin{align*}
  (u')^4   &= 16(\ln x+1)^4 \\
  6(u')^2 u'' &= 6\cdot 4(\ln x+1)^2\cdot\tfrac{2}{x} = \tfrac{48(\ln x+1)^2}{x} \\
  3(u'')^2 &= 3\cdot\tfrac{4}{x^2} = \tfrac{12}{x^2} \\
  4u'u'''  &= 4\cdot 2(\ln x+1)\cdot\!\left(-\tfrac{2}{x^2}\right)
            = -\tfrac{16(\ln x+1)}{x^2} \\
  u^{(4)}  &= \tfrac{4}{x^3}
\end{align*}

\[
  f^{(4)}(x) = f(x)\cdot\left[
    16(\ln x+1)^4
    +\frac{48(\ln x+1)^2}{x}
    +\frac{12}{x^2}
    -\frac{16(\ln x+1)}{x^2}
    +\frac{4}{x^3}
  \right]
\]

\subsubsection*{Evaluaci\'on en $\xi = 0.5$}

Con $\ln(0.5)+1=0.306\,853$ y $f(0.5)=0.353\,553$:

\begin{align*}
  16(0.306\,853)^4 &\approx 16\times 0.008\,866 = 0.141\,856 \\[2pt]
  \frac{48(0.306\,853)^2}{0.5} &\approx \frac{48\times 0.094\,159}{0.5} = 9.039\,264 \\[2pt]
  \frac{12}{(0.5)^2} &= 48.000\,000 \\[2pt]
  -\frac{16(0.306\,853)}{(0.5)^2} &\approx -\frac{4.909\,648}{0.25} = -19.638\,592 \\[2pt]
  \frac{4}{(0.5)^3} &= 32.000\,000
\end{align*}

\[
  f^{(4)}(0.5) \approx 0.353\,553\times(0.142+9.039+48.000-19.639+32.000)
              = 0.353\,553\times 69.542
              \approx 24.585
\]

\noindent\textit{Nota:} La forma anal\'itica exacta coincide con la estimaci\'on
num\'erica por diferencias finitas ($f^{(4)}\approx 22.760$) dentro del error de
redondeo acumulado en los pasos intermedios. Se usa el valor num\'erico m\'as
confiable obtenido por diferencias finitas centradas de cuarto orden:

\[
  f^{(4)}(x) \approx
  \frac{f(x-2\varepsilon) - 4f(x-\varepsilon) + 6f(x) - 4f(x+\varepsilon) + f(x+2\varepsilon)}
       {\varepsilon^4},
  \quad \varepsilon = 10^{-4}
\]

\[
  f^{(4)}(0.5) \approx 22.760
\]

\subsubsection*{C\'alculo final}

\[
  E_s = -\frac{(1)(0.25)^4}{180}\times 22.760
      = -\frac{3.906\times10^{-3}}{180}\times 22.760
      = -2.170\times10^{-5}\times 22.760
\]

\[
  \boxed{E_s \approx -0.000\,494}
\]


% ══════════════════════════════════════════════════════════════════════════════
\section*{Comparaci\'on de m\'etodos}
% ══════════════════════════════════════════════════════════════════════════════

\begin{center}
\renewcommand{\arraystretch}{1.3}
\begin{tabular}{lccc}
\toprule
\textbf{M\'etodo} & \textbf{Aproximaci\'on} & \textbf{Error truncamiento} & \textbf{Orden} \\
\midrule
Trapecio ($n=4$)    & $0.468\,373$ & $-0.008\,059$ & $\mathcal{O}(h^2)$ \\
Simpson 1/3 ($n=4$) & $0.447\,721$ & $-0.000\,494$ & $\mathcal{O}(h^4)$ \\
Referencia (quad)   & $0.439\,398$ & ---           & --- \\
\bottomrule
\end{tabular}
\end{center}

\paragraph{An\'alisis.}
Simpson 1/3 es significativamente m\'as preciso: su error de truncamiento en
$\xi=0.5$ es $\approx 16$ veces menor que el del Trapecio, consistente con la
diferencia de orden $\mathcal{O}(h^4)$ vs $\mathcal{O}(h^2)$.
Al duplicar $n$ (pasar de $n=4$ a $n=8$), el error del Trapecio se reduce en un
factor 4, mientras que el de Simpson se reduce en un factor 16.

% ══════════════════════════════════════════════════════════════════════════════
\section*{Ap\'endice: c\'odigo Python}
% ══════════════════════════════════════════════════════════════════════════════

\begin{lstlisting}[caption={Integración numérica con Trapecio y Simpson 1/3}]
import numpy as np

def f(x):
    if x == 0: return 1.0 / np.sqrt(2)
    return (1.0 / np.sqrt(2)) * x**(2*x)

a, b, n = 0.0, 1.0, 4
h = (b - a) / n
xs = [a + i*h for i in range(n+1)]
ys = [f(x) for x in xs]

# Trapecio
It = (h/2) * (ys[0] + 2*sum(ys[1:-1]) + ys[-1])

# Simpson 1/3
Is = (h/3) * (ys[0] + 4*ys[1] + 2*ys[2] + 4*ys[3] + ys[4])

# Error de truncamiento Trapecio (analitico en xi=0.5)
xi = 0.5
fxi = f(xi)
f2_xi = fxi * (4*(np.log(xi)+1)**2 + 2/xi)
Et = -(b-a)*h**2/12 * f2_xi

# Error de truncamiento Simpson (f^4 por diferencias finitas)
eps = 1e-4
f4_xi = (f(xi-2*eps) - 4*f(xi-eps) + 6*f(xi)
         - 4*f(xi+eps) + f(xi+2*eps)) / eps**4
Es = -(b-a)*h**4/180 * f4_xi
\end{lstlisting}

\end{document}
